% Options for packages loaded elsewhere
\PassOptionsToPackage{unicode}{hyperref}
\PassOptionsToPackage{hyphens}{url}
\documentclass[
]{article}
\usepackage{xcolor}
\usepackage[margin=1in]{geometry}
\usepackage{amsmath,amssymb}
\setcounter{secnumdepth}{-\maxdimen} % remove section numbering
\usepackage{iftex}
\ifPDFTeX
  \usepackage[T1]{fontenc}
  \usepackage[utf8]{inputenc}
  \usepackage{textcomp} % provide euro and other symbols
\else % if luatex or xetex
  \usepackage{unicode-math} % this also loads fontspec
  \defaultfontfeatures{Scale=MatchLowercase}
  \defaultfontfeatures[\rmfamily]{Ligatures=TeX,Scale=1}
\fi
\usepackage{lmodern}
\ifPDFTeX\else
  % xetex/luatex font selection
\fi
% Use upquote if available, for straight quotes in verbatim environments
\IfFileExists{upquote.sty}{\usepackage{upquote}}{}
\IfFileExists{microtype.sty}{% use microtype if available
  \usepackage[]{microtype}
  \UseMicrotypeSet[protrusion]{basicmath} % disable protrusion for tt fonts
}{}
\makeatletter
\@ifundefined{KOMAClassName}{% if non-KOMA class
  \IfFileExists{parskip.sty}{%
    \usepackage{parskip}
  }{% else
    \setlength{\parindent}{0pt}
    \setlength{\parskip}{6pt plus 2pt minus 1pt}}
}{% if KOMA class
  \KOMAoptions{parskip=half}}
\makeatother
\usepackage{longtable,booktabs,array}
\newcounter{none} % for unnumbered tables
\usepackage{calc} % for calculating minipage widths
% Correct order of tables after \paragraph or \subparagraph
\usepackage{etoolbox}
\makeatletter
\patchcmd\longtable{\par}{\if@noskipsec\mbox{}\fi\par}{}{}
\makeatother
% Allow footnotes in longtable head/foot
\IfFileExists{footnotehyper.sty}{\usepackage{footnotehyper}}{\usepackage{footnote}}
\makesavenoteenv{longtable}
\setlength{\emergencystretch}{3em} % prevent overfull lines
\providecommand{\tightlist}{%
  \setlength{\itemsep}{0pt}\setlength{\parskip}{0pt}}
\usepackage{bookmark}
\IfFileExists{xurl.sty}{\usepackage{xurl}}{} % add URL line breaks if available
\urlstyle{same}
\hypersetup{
  pdftitle={An AI oracle in the style of the generic group model},
  hidelinks,
  pdfcreator={LaTeX via pandoc}}

\title{An AI oracle in the style of the generic group model}
\author{}
\date{2026-10-13}

\begin{document}
\maketitle

\section{Why the generic group model and not the random oracle
model}\label{why-the-generic-group-model-and-not-the-random-oracle-model}

The random oracle model idealizes the \emph{values} of a hash function:
they are uniform and independent. Trained models have no such randomness
to idealize. The generic group model (Shoup, EUROCRYPT 1997; Maurer, IMA
2005; for rings, Aggarwal and Maurer, EUROCRYPT 2009) idealizes
\emph{access}. Every party holds opaque handles to group elements and
may only apply the group operations through the oracle. Lower bounds
become provable, and they hold only for algorithms that never look at
the representation. Index calculus, which does look at the
representation of \(\mathbb{Z}_p^*\), beats the generic bound. That is
the black-box against white-box question in another setting, and the
reason this note copies the generic model.

\section{Objects}\label{objects}

\begin{itemize}
\tightlist
\item
  A data distribution \(\mathcal{D}\) over labelled points, and a
  training set \(S\) of \(n\) points.
\item
  An architecture: a parameter space \(\Theta\) and a map
  \(\theta \mapsto f_\theta\) from parameters to functions. Write
  \(\mathcal{F} = \{ f_\theta : \theta \in \Theta \}\).
\item
  The \emph{fibre} of a function \(f \in \mathcal{F}\) is
  \(\Phi(f) = \{ \theta \in \Theta : f_\theta = f \}\), all parameter
  vectors that compute it. This is the fibre of the parameterization,
  not of the training map.
\item
  A training procedure \(\mathsf{Train}\), a randomized map from
  training sets to \(\Theta\). It is never opened. It is constrained
  only by utility.
\end{itemize}

\textbf{Utility condition \(\alpha\).}
\(\mathbb{E}_{S \sim \mathcal{D}^n}\big[\Pr_{z \sim \mathcal{D}}[f_{\mathsf{Train}(S)} \text{ correct on } z]\big] \ge \alpha\).
Without it a constant output is perfectly private and nothing is
learned.

\section{The oracle}\label{the-oracle}

The oracle \(\mathcal{O}\) keeps a private table from handles to
parameter vectors and offers these operations.

{\def\LTcaptype{none} % do not increment counter
\begin{longtable}[]{@{}
  >{\raggedright\arraybackslash}p{(\linewidth - 4\tabcolsep) * \real{0.3333}}
  >{\raggedright\arraybackslash}p{(\linewidth - 4\tabcolsep) * \real{0.3333}}
  >{\raggedright\arraybackslash}p{(\linewidth - 4\tabcolsep) * \real{0.3333}}@{}}
\toprule\noalign{}
\begin{minipage}[b]{\linewidth}\raggedright
Operation
\end{minipage} & \begin{minipage}[b]{\linewidth}\raggedright
Returns
\end{minipage} & \begin{minipage}[b]{\linewidth}\raggedright
Level
\end{minipage} \\
\midrule\noalign{}
\endhead
\bottomrule\noalign{}
\endlastfoot
\(\mathsf{train}(S)\) & a fresh handle to
\(\theta \leftarrow \mathsf{Train}(S)\) & call \\
\(\mathsf{eval}(h, x)\) & \(f_\theta(x)\) & function \\
\(\mathsf{grad}_x(h, z)\) & gradient of the loss in the input &
function \\
\(\mathsf{grad}_\theta(h, z)\) & gradient of the loss in the parameters
& representation \\
\(\mathsf{reveal}(h)\) & \(\theta\) & representation \\
\end{longtable}
}

An \emph{access class} is a set of permitted operations with a budget on
each. Black-box is \(\{\mathsf{eval}\}\). The gradient attacks of the
literature use \(\{\mathsf{eval}, \mathsf{grad}_\theta\}\) or
\(\{\mathsf{eval}, \mathsf{grad}_x\}\), and the two are different
classes. Shadow-model attacks add \(\mathsf{train}\) on chosen sets.
White-box is \(\{\mathsf{reveal}\}\).

\section{The game}\label{the-game}

The challenger draws \(S \sim \mathcal{D}^n\), \(z_1\) uniformly from
\(S\), \(z_0 \sim \mathcal{D}\), and \(b\) uniformly. It gives the
adversary \(z_b\) and a handle to \(\mathsf{Train}(S)\). The adversary
uses its access class and outputs \(b'\). Its advantage is
\(2\Pr[b'=b]-1\).

\section{Ideal representation}\label{ideal-representation}

\textbf{Assumption.} Conditioned on the function
\(f_{\mathsf{Train}(S)}\), the parameter vector \(\mathsf{Train}(S)\) is
independent of \(S\). One instance: it is uniform on the fibre
\(\Phi(f_{\mathsf{Train}(S)})\).

This is the idealization, in the place the random encoding holds in the
generic group model. It says the weights carry nothing about the data
beyond the function they compute.

\textbf{Proposition.} Under ideal representation, the best unbounded
white-box advantage equals the best unbounded advantage of an attacker
that sees only the function \(f_{\mathsf{Train}(S)}\).

\emph{Proof.} Under the assumption, the membership bit, the function and
the parameters form a Markov chain. The best advantage of an unbounded
attacker equals the total variation distance between the two conditional
distributions of its view. Total variation does not increase under a
channel, and the parameters are a channel applied to the function. The
white-box view therefore gives at most the function's advantage, and it
gives at least that much because the function is computable from the
parameters. \(\square\)

\section{What the definition
separates}\label{what-the-definition-separates}

Real training breaks the assumption, and the white-box against black-box
gap splits into two parts.

\begin{enumerate}
\def\labelenumi{\arabic{enumi}.}
\tightlist
\item
  \textbf{Representation leakage.} What the parameters say about
  membership beyond the function. It is zero when the assumption holds,
  and it is the information gap.
\item
  \textbf{Extraction cost.} With representation leakage zero, an
  efficient black-box attacker may still fail to compute from queries
  what the parameters hand over for free. This is the computational gap.
  For ReLU networks, polynomial-time extraction up to the fibre (CJM20,
  CCHRSS24) suggests it is small. For function families that resist
  learning from queries while their code reveals a secret (BGIRSVY01) it
  can be large.
\end{enumerate}

\section{The methodology caveat}\label{the-methodology-caveat}

An oracle that hands out an ideal representation of every trained
function is a virtual black-box obfuscator for the trained function.
BGIRSVY01 proves such obfuscators do not exist for all function
families. So the ideal representation cannot be instantiated in general,
as the random oracle cannot (Canetti, Goldreich, Halevi, STOC 1998).
This is an analogy, recorded as heuristic, and it says the definition
will have the same standing as the random oracle model: a model in which
proofs are possible, with known families where it fails.

\section{Open}\label{open}

\begin{itemize}
\tightlist
\item
  Which access classes are strictly ordered for natural
  \((\mathcal{D}, \Theta, \alpha)\), in particular
  \(\{\mathsf{eval}, \mathsf{grad}_x\}\) against
  \(\{\mathsf{eval}, \mathsf{grad}_\theta\}\).
\item
  Lower bounds on \(\mathsf{eval}\) and \(\mathsf{train}\) calls for
  membership inference inside the model, in the way the generic group
  model gives a square-root lower bound for discrete logarithms.
\item
  Whether any natural training procedure satisfies ideal representation,
  even approximately.
\end{itemize}

\end{document}
